\documentclass{situationist-preprint}

\title{Metaphysical Indeterminacy, Restricted}
\msid{MET-19-a}
\msfamily{MET-19}
\msdiscipline{Metaphysics}
\mstype{Principal manuscript}
\mskeywords{metaphysical indeterminacy; states of affairs; self-reference; groundedness; liar paradox; vagueness; open future}

\begin{document}
\maketitle

\begin{abstract}
Is there indeterminacy in the world itself, and not only in our representations of it? We argue that there is, but only in a restricted class of cases: those in which whether a state of affairs obtains depends on whether that very state of affairs obtains. A state of affairs such as that every state of affairs Jones is now considering fails to obtain, when it is the only one Jones is considering, can neither obtain nor fail to obtain on pain of contradiction, and since there are no true contradictions it is neither. The indeterminacy is neither linguistic nor epistemic. We then argue that no other source of indeterminacy is metaphysical: the indeterminacy of vagueness is epistemic, the openness of the future is unsettledness at a moment, and quantum mechanics on the Everett interpretation describes a determinate world. Indeterminacy arises only where obtaining is not grounded in states of affairs other than the one in question.
\end{abstract}

\section{The question}

Metaphysical indeterminacy is indeterminacy not due to the way the world is represented or known. \citet{BarnesWilliams2011} and \citet{Wilson2013} have given theories of it, and its proposed instances include vague objects, the open future, quantum indeterminacy and unsettled matters in fiction and law. The opposing view, that the world is fully determinate and all indeterminacy lies in representation, goes back to \citet{Russell1923} and is defended for the case of vagueness in MET-18. We argue for an intermediate position: metaphysical indeterminacy is real, but its only source is self-reference among states of affairs.

\section{Framework}

\begin{definition}\label{d:soa}
A \emph{state of affairs} is a way for things to be, which obtains or fails to obtain \citep{Armstrong1997}. A state of affairs $S$ is \emph{self-referential} if what is required for $S$ to obtain includes a condition on whether $S$ itself obtains.
\end{definition}

\begin{principle}[Obtaining]\label{p:obt}
For any state of affairs $S$ that things are thus-and-so, $S$ obtains if and only if things are thus-and-so.
\end{principle}

\begin{principle}[Exclusion]\label{p:excl}
No state of affairs both obtains and fails to obtain.
\end{principle}

\ref{p:obt} is the worldly counterpart of the truth schema; it is constitutive of the notion of a state of affairs. \ref{p:excl} is the worldly counterpart of the principle defended in LOG-03, on which there are no true contradictions.

\section{The argument}

\begin{proposition}[Existence]\label{p:exist}
There are self-referential states of affairs.
\end{proposition}

\begin{proof}
Self-reference among states of affairs need not be engineered by language; as \citet{Kripke1975} observed for sentences, it can arise from contingent facts. Suppose Jones is considering, at $t$, exactly one state of affairs, and that it is $L$: that every state of affairs Jones is considering at $t$ fails to obtain. Jones's considering a state of affairs is a mental relation to it, not a sentence. $L$ exists if states of affairs are as abundant as the ways things can be, and what is required for $L$ to obtain is a condition on whether $L$ obtains.
\end{proof}

\begin{theorem}\label{t:main}
$L$ neither obtains nor fails to obtain, and this indeterminacy is metaphysical.
\end{theorem}

\begin{proof}
By \ref{p:obt}, $L$ obtains if and only if every state of affairs Jones is considering at $t$ fails to obtain, that is, if and only if $L$ fails to obtain. If $L$ obtained, it would fail to obtain; if it failed to obtain, it would obtain; either way it would both obtain and fail to obtain, contrary to \ref{p:excl}. So $L$ neither obtains nor fails to obtain. The indeterminacy is not linguistic, since $L$ is not a sentence and the argument uses no semantic predicate, only the worldly predicate \enquote{obtains}. It is not epistemic, since no fact about $L$ is unknown.
\end{proof}

Theorem~\ref{t:main} is the worldly analogue of the treatment of the liar in LOG-01-a, on which the liar sentence has an intermediate value in Łukasiewicz logic, strictly between truth and falsity, and revenge is blocked by a hierarchy of determinateness operators. We take the intermediate value to be the status of $L$: it is not determinately obtaining and not determinately failing to obtain. The case is consistent with LNG-18, on which self-reference produces paradox only together with a semantic predicate; \enquote{obtains} plays here the role that \enquote{true} plays for sentences.

\begin{principle}[Grounded determinacy]\label{p:ground}
If whether $S$ obtains is fixed by whether certain other states of affairs obtain, none of which is $S$ and each of which determinately obtains or fails to obtain, then $S$ determinately obtains or fails to obtain.
\end{principle}

\ref{p:ground} is the worldly form of Kripke's notion of groundedness \citep{Kripke1975}: determinacy is inherited along non-circular dependence. It implies that indeterminacy can enter only at states of affairs that are not fixed by others, or at states of affairs whose determinants include themselves. The first are the basic states of affairs of things having properties and standing in relations. We now argue that none of the proposed sources makes any of these indeterminate.

\begin{proposition}\label{p:restrict}
Vagueness, the open future and quantum mechanics do not produce metaphysical indeterminacy.
\end{proposition}

\begin{proof}[Argument]
Vagueness: by MET-18, vague predicates have sharp boundaries that cannot be known, and there are no vague objects; the indeterminacy of borderline cases is epistemic. The open future: ACT-17 and TSM-12 hold that future contingents lack a truth value before they are settled. As REL-05-a reads ACT-17 in branching-time semantics, a future-tensed sentence is evaluated at a moment relative to the histories through it, and lacks a value at the moment when it is true on some and false on others; a tenseless proposition is true relative to the actual history. What lacks a value is the future-tensed sentence evaluated at a moment, not a state of affairs, located at a time, that neither obtains nor fails to obtain. This is unsettledness at a moment, and it is a feature of tensed evaluation. Quantum mechanics: on the Everett interpretation defended in SCI-09, the quantum state evolves deterministically and every branch is determinate. Unsettled matters in fiction and law concern what a body of representations settles, not the world.
\end{proof}

By \ref{p:ground} and Proposition~\ref{p:restrict}, metaphysical indeterminacy is confined to states of affairs whose determinants include themselves. These are exactly the self-referential states of affairs of Definition~\ref{d:soa}, and Theorem~\ref{t:main} shows that some of them are indeterminate.

\section{Objections}

\begin{objection}
If there is any metaphysical indeterminacy, Evans's argument \citep{Evans1978} threatens: indeterminacy in the world should be capable of applying to identity, which MET-18 shows incoherent.
\end{objection}

\begin{reply}
Evans's argument shows that an identity cannot be indeterminate where that indeterminacy is determinate. The indeterminacy defended here concerns whether certain states of affairs obtain, and arises from their self-reference; no identity statement among objects is self-referential in this way. Nothing here implies vague objects, and MET-18's conclusion stands.
\end{reply}

\begin{objection}
Self-referential states of affairs should be denied existence, as some deny that the liar sentence expresses a proposition.
\end{objection}

\begin{reply}
Whether $L$ is self-referential depends on what Jones happens to consider. A restriction that excluded it would exclude innocent states of affairs that differ from it only in that accident, and so give up the abundance on which the theory of states of affairs relies.
\end{reply}

\section{Conclusion}

There is metaphysical indeterminacy. It arises when a state of affairs has its own obtaining among its conditions, as in the worldly analogue of the liar, and in no other case: the indeterminacy of vagueness is epistemic, the openness of the future is a feature of tensed evaluation, and the quantum world is determinate. Grounded states of affairs inherit determinacy, and only the ungrounded, self-referential ones can lack it.

\begin{thebibliography}{99}\small
\bibitem[Armstrong(1997)]{Armstrong1997} Armstrong, D.~M. (1997). \emph{A World of States of Affairs}. Cambridge University Press.
\bibitem[Barnes and Williams(2011)]{BarnesWilliams2011} Barnes, E. and Williams, J.~R.~G. (2011). A theory of metaphysical indeterminacy. \emph{Oxford Studies in Metaphysics}, 6, 103--148.
\bibitem[Evans(1978)]{Evans1978} Evans, G. (1978). Can there be vague objects? \emph{Analysis}, 38(4), 208.
\bibitem[Kripke(1975)]{Kripke1975} Kripke, S. (1975). Outline of a theory of truth. \emph{The Journal of Philosophy}, 72(19), 690--716.
\bibitem[Russell(1923)]{Russell1923} Russell, B. (1923). Vagueness. \emph{Australasian Journal of Psychology and Philosophy}, 1(2), 84--92.
\bibitem[Wilson(2013)]{Wilson2013} Wilson, J.~M. (2013). A determinable-based account of metaphysical indeterminacy. \emph{Inquiry}, 56(4), 359--385.
\end{thebibliography}

\end{document}
